\subsection{One source-bound boundary query with endpoint moments}
\label{marked-cut:moments}

Report 93 supplies a native adapter for moment-coded first returns: it starts
with the actual parallelity polygons and free-edge records, rather than an
abstract path pairing. We integrate its producer, independent checker,
literal path oracle and arithmetic tests. The supplied-cut operator exposes
\path{boundary_method='moments'} alongside the retained mask method. The
new compiler also records the exact number of original-boundary edges in
every gap. This is local length and fibre data, not transport of an arbitrary
pattern hidden inside a gap.

\paragraph{The source boundary packet.}
Let $p$ be the number of core-facing ports and $N$ the corner-occurrence
namespace size. Both are determined by the validated source triangulation
and binary normal coordinates. The exceptional port list has size $O(T)$;
$N$ itself may be exponentially larger than the input bit length. All source
gluings are interval maps with separate endpoint-reversal and fibre-flip
flags. In particular, fibre transport cannot be inferred from the slope of
a singleton interval map.

Delete the open port edges from the boundary of every parallelity base.
Each active residual component is a path with two endpoint occurrences;
a single point is a permitted path. Components without endpoints are passive
circles or silent occurrence classes. Work on the two-sheeted fibre cover,
so the active paths have $r=4p$ distinct endpoint labels, numbered $1$ through
$r$. Each active path downstairs has two compatible lifts. The source
surface and edge-link checks, not a numerical moment identity, establish
this geometric degree-two property.

\begin{proposition}[Canonical source-edge mass]
Each original-boundary edge occurrence can be assigned exactly once to its
local initial corner. This assignment is represented by $O(T)$ intervals
in the corner namespace, without enumerating their binary lengths.
\end{proposition}
\begin{proof}
For an original-boundary record on edge $j$ of polygon family $P$, a rank
$k$ is owned at $c_{P,j}+k-1$, where $c_{P,j}$ is that corner slot's start.
A rank interval $[u,v)$ therefore contributes the owner interval
$[c_{P,j}+u-1,c_{P,j}+v-1)$. Source free-edge records are disjoint rank gaps;
different polygon corner slots occupy disjoint namespaces. Thus every edge
has one owner and no owner is counted twice before quotienting. Glued
corners can identify several owners, and then their masses add, as required.
The same outgoing-corner convention also assigns the port edges disjoint
owner occurrences. If $E$ is the number of represented original-boundary
edges, then $E+p\le N$. Each lift inherits one copy of this ownership.
\end{proof}

\paragraph{The pair codec.}
For $p>0$, choose the power-of-two radices
\[
 A=2^{\operatorname{bitlength}(2r)},\qquad
 s_i=i+A i^2,\qquad
 D=2^{\operatorname{bitlength}(2(r+A r^2))}.
\]
Label $i$ contributes $s_i$ at its source endpoint occurrence. Each lifted
owner edge contributes $D$. Run one scalar weighted orbit census on the
residual double cover. A path of length $\ell$ with endpoint labels $i<j$
has weight $W=\ell D+s_i+s_j$.

\begin{proposition}[Exact pair and length decoding]
For an authenticated active path, its two endpoint labels and exact edge
length are recovered by
\[
 \ell=\lfloor W/D\rfloor,\quad R=W\bmod D,\quad
 S=R\bmod A,\quad Q=\lfloor R/A\rfloor,\quad
 i,j=\frac{S\mp\sqrt{2Q-S^2}}2.
\]
All calculations use exact integer arithmetic. Passive components have
remainder zero and do not contribute a port pairing.
\end{proposition}
\begin{proof}
We have $i+j\le2r<A$ and $s_i+s_j\le2(r+A r^2)<D$. Thus neither radix has a
carry from the two endpoint contributions. The remainder determines
$S=i+j$ and $Q=i^2+j^2$, and $2Q-S^2=(j-i)^2$. Distinct positive labels give
a positive square discriminant. A component without labels has weight
$\ell D$ and remainder zero. Canonical ownership makes $\ell$ the count of
original edges, rather than the count of aliased corner occurrences.
\end{proof}

The verifier requires a positive square discriminant, parity, label range,
active-row multiplicity one, and unique complete coverage of all labels.
It checks equal lengths and compatible images for the two lifts of each
endpoint, before traversing the boundary permutation independently. It
uses the product $(S^2-Q)/2$ and the roots of $x^2-Sx+(S^2-Q)/2$, rather than
calling the producer's shift/mask decoder. No floating-point square roots
or large endpoint masks are used.

\begin{proposition}[Why the global inventory check is essential]
Consider any partition of the actual unique labels and any nonnegative edge
masses. If the pair decoder accepts a complete disjoint matching of all $r$
labels, every accepted active row contains exactly two actual labels, and
every ignored row contains none.
\end{proposition}
\begin{proof}
An actual zero-label row has remainder zero and cannot supply a decoded
pair. An actual singleton has remainder $s_i$ without a carry; its formal
roots are $0,i$, which fail the positive-label check. Hence every accepted
active row consumes at least two actual labels. The complete disjoint
decoded matching has $r/2$ rows, requiring at least all $r$ actual labels.
Since actual rows partition those labels, every active row has exactly two
and no labels can remain in ignored rows. The preceding no-carry proof now
applies to each active row. Without complete unique coverage, a partial
moment record would not justify this counting argument.
\end{proof}

\paragraph{What replaces the corner query.}
A gap is a point exactly when $\ell=0$. Thus positive length gives precisely
the legacy \path{old_boundary_present} flag. Matching endpoints, following
a port to its other endpoint, and then crossing the paired gap reconstruct
cyclic words and maximal attaching intervals. Local fibre flips are retained;
the total flip around an authenticated boundary circle is zero. Removing
only \path{unmarked_edges} from each gap gives the complete previous
skeleton. In particular, the moment compiler does not infer attachment
order from a histogram of lengths alone.

If $p=0$, the active skeleton is empty and no boundary census is needed.
The standalone packet adapter additionally binds active loops to the
already certified base profiles and retains passive boundary multiplicities.
It authenticates the four upstream source traces; that work is additional
to the claim of one new boundary census.

\paragraph{Bit size and the remaining scheduling obligation.}
The pair scalar has $D=O(p^3)$. Total initial mass is
$2ED+\sum_{i=1}^r s_i=O(Np^3)$, using $p\le N$. Thus nonnegative
mass-preserving operations need $O(\log(N+1)+\log(p+1))$ bits for scalar
weights. This replaces the linear-in-$p$ endpoint-mask width; it also adds
exact edge lengths that the old scalar did not contain. The retained
four-field count/sum/squares/length codec has the larger magnitude bound
$O(Np^9)$ and serves as an independent differential reference.

Query construction uses $O(T)$ interval records and $O(p)$ endpoint records.
Decoding and word traversal use polynomial bit work and linear structural
work in their records. However, ownership intervals can increase the number
of intermediate weight runs. The old sparse point-tag bounds do not apply
to this dense mass program. The weighted orbit theorem~\cite{aht-orbits} supplies a polynomial
algorithm for the encoded task, but the present integration does not prove
a uniform polynomial bound for every maintained scheduler trace. Actual
cycle counts, weight runs and certificate sizes must be reported separately.

\paragraph{Native integration and proof budgets.}
The whole-cut moment path reuses its own already constructed parallelity
geometry and source fingerprint. It avoids reconstructing that geometry
in a second producer or replaying an external packet that it did not receive.
The independent checker still reconstructs the source geometry and checks
the weighted proof. New nested boundary proofs have schema
\path{native-moment-boundary-v2}; the outer v1--v6 cut formats and their
geometric meanings remain unchanged. Existing mask proofs remain supported.

The four source queries and the one moment query share the whole-cut cycle
allowance. The replay cap counts all five traces, including the weighted
traces inside the two source selectors. Port-free boundary proofs contribute
zero operations. Exhaustion returns no partial cut certificate; invalid
backend names and callback interruption do not silently choose another path.

\paragraph{Validation and measurements.}
The nine delivered moment tests pass against the maintained runtime in
159.509 seconds from an isolated copy, leaving the delivered report files
unchanged. The ported nine native tests pass in 164.451 seconds, and five
whole-cut integration tests pass in 6.662 seconds. Coverage includes all
157 inherited boundary sources, both scalar codecs, 146 literal path
comparisons, binary multiplicity $2^{16384}$, hostile proofs, independent
replay with producers disabled, and the shared five-trace budget.
The arithmetic test checks all 4,140 partitions of eight endpoint labels:
exactly the 105 perfect matchings are accepted.
The full maintained suite passes all 1,710 tests in 735.091 seconds, with
no skips. It also reports 13,846 expanded-sheet topology comparisons and
196,712 independent marked-cover comparisons.

The complete differential audit checks 106 historical sources and 14
interior-vertex fixtures. Both backends complete the same 102 cuts; all
complete censuses agree, all certificates replay, and removing just the new
gap-length field recovers every old boundary skeleton. All 92 historical
complete mask certificates are byte-identical. The 95 completed cases with
retained literal Regina references agree with those references; seven
large-coordinate cases have no expanded reference. This comparison uses
the previously recorded oracle outputs on hash-identical source files,
not a fresh Regina run. The remaining 18 cases stay inconclusive.

\paragraph{Complete-operation timings.}
The baseline consists of the producer and verifier from commit
\texttt{a1f5adb096fa}, with their unchanged lower-level dependencies.
Each operation produces the full cut certificate and independently replays
it. There are five randomly ordered pairs and three mask/mask controls per
source, in one process, before starting the full suite or corpus audit.
The seven families were selected in advance; the cap-only representative
is the first semantically eligible completed source, not a timing selection.
The moment path combines the one-query encoding with reuse of geometry
already built by its producer. These measurements do not isolate the two
effects.

\begin{center}
\small
\begin{tabular}{lrrrrr}
\toprule
Source & Ports & Masks (ms) & Moments (ms) & Paired gain & A/A \\
\midrule
Layered pilot & 16 & 40.25 & 35.52 & 1.142 & 1.012 \\
Fibonacci 9, mixed & 148 & 2351.94 & 2165.22 & 1.087 & 1.019 \\
Interior, mixed & 11 & 201.93 & 192.68 & 0.947 & 0.994 \\
Figure-eight, mixed & 194 & 6214.97 & 5572.34 & 1.119 & 0.975 \\
Layered, $2^{512}$ & 2 & 125.15 & 139.84 & 0.932 & 0.973 \\
Port-free & 0 & 2.16 & 1.93 & 1.122 & 0.995 \\
Cap-only & 8 & 7.85 & 7.27 & 1.079 & 0.985 \\
\bottomrule
\end{tabular}
\end{center}


The reported gain is the median of the five paired mask/moment ratios,
not the ratio of the two marginal medians. In particular, the interior
example has a paired ratio below one despite its lower marginal moment
median. Its samples span 182--352 ms for masks and 188--380 ms for moments.
The figure-eight ranges are 6.16--7.21 s and 5.49--7.10 s. The port-free
moment case has one 8.91 ms outlier, despite a 1.93 ms median. The complete
raw samples and execution order are retained. Five paired medians favor
moments, while two favor masks; the mask/mask control medians range from
0.973 to 1.019. Accordingly \path{boundary_method='masks'} remains the
default and \path{'moments'} is an explicit option.

\paragraph{Structural cost and certificate size.}
For the figure-eight case, the two old boundary queries total 1,080 orbit
cycles; the new one uses 733. Its maximum weight-run count rises from 881
in the old cover query to 919 in the new one. The binary layered example
has 168 versus 131 cycles, but cover weight runs rise from 17 to 85.
This explains why fewer queries alone need not yield lower elapsed time.
All seven full certificates shrink, by 0.39--19.65 percent under the same
JSON serialization. The figure-eight proof drops from 1,070,683 to 860,263
bytes, and the binary example from 490,009 to 460,255 bytes.

For the figure-eight input, 776 endpoint labels use a maximum additive
source weight of 33 bits. This statistic is not the maximum transported
mass: the proved total-mass bound above supplies that distinction.
The delivered binary test, rerun against the maintained orbit engine,
uses 131 cycles, 60 owner intervals and 68 additive weight records for
each of the multiplicities $2^{32}$, $2^{512}$ and $2^{16384}$. Those
three observations do not establish a uniform scheduler bound.

The audit, paired driver, source pins, retained baseline, literal coverage
and test logs are in \path{synthesis/data/moment-boundary-*}. These are
supplied-cut operations; no new surface-discovery algorithm or general
quasi-polynomial recognition bound follows.

